% =============================================================================
%  ¿De dónde salen los futbolistas argentinos?
%  Compilar con:  pdflatex paper.tex   (dos veces, por las referencias cruzadas)
%  Las figuras se toman de ../outputs/figures/ en PDF vectorial.
% =============================================================================
\documentclass[11pt,a4paper]{article}

\usepackage[T1]{fontenc}
\usepackage[utf8]{inputenc}
\usepackage[spanish,es-nodecimaldot,es-tabla]{babel}
\usepackage{lmodern}
\usepackage[a4paper,margin=2.4cm]{geometry}
\usepackage{graphicx}
\usepackage{booktabs}
\usepackage{siunitx}
\usepackage{caption}
\usepackage[hidelinks]{hyperref}
\usepackage{xcolor}
\usepackage{microtype}
\usepackage{enumitem}

\graphicspath{{../outputs/figures/}}

% Las figuras ya traen su propio título y su pie dibujados adentro del PDF: el
% caption de LaTeX solo tiene que numerarlas y decir de dónde salen, sin repetir.
\captionsetup{font=small,labelfont=bf,skip=6pt}
\setlength{\parskip}{0.35em}
\sisetup{output-decimal-marker={,},group-separator={.},group-minimum-digits=4}

\newcommand{\ic}[2]{\mbox{IC 95\,\% #1--#2}}

\title{\bfseries ¿De dónde salen los futbolistas argentinos?\\[0.35em]
\large Geografía del nacimiento y de la formación en el fútbol argentino,\\
cohortes 1975--2008}
\author{Agustín Viullet}
\date{1 de agosto de 2026\\[0.4em]
\footnotesize Análisis reproducible sobre Wikidata, la serie histórica de nacidos
vivos del DEIS\\ y el Censo Nacional 2022 del INDEC. Snapshot de Wikidata:
30 de julio de 2026.}

\begin{document}
\maketitle

\begin{abstract}
\noindent
Se analiza el lugar de nacimiento de \num{5511} futbolistas profesionales
argentinos nacidos entre 1975 y 2008, contra el número de \textbf{nacidos vivos}
de cada cohorte en cada lugar. La tasa se lee directo: de cada \num{100000} bebés
nacidos en un lugar, cuántos llegaron a futbolistas profesionales.

\textbf{El \emph{birthplace effect} clásico no aparece en Argentina y lo que
aparece es su inverso}, pero con dos matices que la literatura previa no
registra. La producción va de 12,9 futbolistas cada \num{100000} nacidos en
localidades de menos de \num{10000} habitantes a 30,5 en aglomerados de más de
\num{500000} (RR 0,42; \ic{0,38}{0,47}); no hay pico en las ciudades medianas.
Sin embargo el efecto \textbf{no es un gradiente sino un escalón} —los nueve
deciles de tamaño por debajo de \num{10000} habitantes no muestran tendencia
alguna— y el tamaño de la ciudad explica apenas el 1\,\% de la variación entre
ciudades (pseudo-$R^2 = 0{,}011$).

\textbf{El resultado más robusto del trabajo es condicional y no usa denominador
poblacional}: entre los futbolistas que ya llegaron a un juvenil de la selección,
los nacidos fuera de un gran aglomerado llegan a la Mayor con más frecuencia
—41,9\,\% contra 28,1\,\%; OR 1,85 (\ic{1,14}{2,98}), $p = 0{,}013$, y 1,85
ajustando por cohorte de nacimiento—. Al jugador del interior le cuesta mucho más
entrar, pero el que entra rinde más. Ese diseño no lo afectan ni el sesgo de
imputación de nacimientos, ni el registro del parto en la ciudad cabecera, ni la
cobertura del corpus.

Se discute en detalle por qué las tasas por nacido —y no el resultado
condicional— quedan expuestas a un artefacto del sistema de salud, y qué haría
falta para descartarlo.
\end{abstract}

\vspace{0.5em}
\hrule
\vspace{1.2em}

% =============================================================================
\section{Introducción}

\subsection{El \emph{birthplace effect}}

El \emph{birthplace effect} —también llamado \emph{place of early development
effect}— es uno de los hallazgos más replicados en la investigación sobre
desarrollo deportivo. Côté y colegas \cite{cote2006} documentaron que las
ciudades chicas y medianas producen desproporcionadamente más atletas de elite
que las grandes urbes y que las zonas rurales, con un óptimo típicamente situado
entre los \num{50000} y los \num{100000} habitantes. La explicación habitual
combina espacio físico para el juego libre, densidad social suficiente para
sostener competencia organizada sin generar barreras de acceso, y relaciones
entrenador-jugador más estables.

La evidencia específica en fútbol es, sin embargo, menos unánime de lo que suele
citarse. La revisión sistemática de Hernández-Simal y colegas \cite{revision2024}
identifica tres ejes —tamaño y densidad poblacional, características
sociodemográficas y \textbf{proximidad a centros de rendimiento}— y reporta
hallazgos contradictorios: hay estudios que encuentran jugadores de elite
proviniendo de áreas densamente pobladas, y concluye que en fútbol no existe una
ventaja consistente de las ciudades chicas.

\subsection{El caso argentino}

Argentina exporta futbolistas en volumen y tiene un relato nacional muy
establecido sobre su origen. El «crack del interior», el potrero, el pueblo chico
que da campeones del mundo —Gualeguay, \num{44000} habitantes, dos— describe con
precisión el patrón que el \emph{birthplace effect} predice.

Ese relato nunca fue puesto a prueba. Lo que existe es periodismo descriptivo:
listas de jugadores por provincia, sin denominador y por lo tanto sin capacidad
de distinguir «produce muchos» de «nace mucha gente». Argentina no aparece en la
revisión sistemática citada. Este trabajo es, hasta donde alcanza nuestra
búsqueda, el primer análisis estadístico del \emph{birthplace effect} en el
fútbol argentino.

\subsection{Hipótesis}

\begin{description}[leftmargin=2.2em,style=nextline,itemsep=0.15em]
  \item[H1] Los futbolistas están sobrerrepresentados entre los nacidos en
    ciudades chicas y medianas.
  \item[H2] El interior produce más futbolistas per cápita que el AMBA.
  \item[H3] Existe migración sistemática entre el lugar de nacimiento y el club
    formador.
  \item[H4] El efecto se intensifica con el nivel competitivo alcanzado.
  \item[Exploratorio] Asociación entre posición y región de nacimiento.
\end{description}

% =============================================================================
\section{Métodos}

\subsection{El denominador: nacidos vivos, no población censada}
\label{sec:denominador}

Es la decisión metodológica central y conviene explicar por qué.

\textbf{Lo que no sirve.} Contar futbolistas nacidos en 1975 contra la población
censada en 2022 no mide nacimientos: mide quiénes seguían vivos y residiendo en
el mismo lugar casi cincuenta años después. Ese denominador está contaminado por
mortalidad y, sobre todo, por migración interna, que en Argentina va justamente
del interior hacia el centro, en la misma dirección que el fenómeno bajo estudio.

\textbf{Lo que se usa.} El denominador es el número de nacidos vivos de cada
cohorte en cada lugar:

\begin{center}\small
\begin{tabular}{@{}lp{7.6cm}l@{}}
\toprule
\textbf{Nivel} & \textbf{Fuente} & \textbf{Naturaleza}\\
\midrule
Provincia & DEIS, serie histórica de nacidos vivos por jurisdicción, 1914--2024
  & \textbf{Dato real}\\[0.4em]
Departamento / ciudad & el total provincial real, repartido según la
  participación de cada departamento en la población de su provincia en el censo
  más cercano al año de nacimiento (1991, 2001, 2010, 2022) & \emph{Estimado}\\
\bottomrule
\end{tabular}
\end{center}

El reparto intraprovincial es el único supuesto de la cadena y se valida contra
los nacimientos departamentales reales del RENAPER (2012--2022): el error
relativo mediano es del 9,1\,\% y el 83,8\,\% de los departamentos cae dentro del
20\,\%.

\textbf{Ese error tiene signo y tiene pendiente, y hay que decirlo}
(figura~\ref{fig:sesgo}). No se distribuye parejo: al decil de departamentos más
chicos el estimador le asigna un \textbf{17\,\% más} de nacimientos de los
reales, mientras que en el decil más grande no se equivoca (Spearman entre tamaño
y ratio estimado/real $= -0{,}355$; $p < 10^{-15}$). Un denominador inflado
deprime la tasa, de modo que \textbf{el sesgo empuja en la misma dirección que el
hallazgo principal}: corregirlo llevaría el RR de las localidades chicas de 0,42
a aproximadamente 0,49. El efecto sobreviviría, pero es un 17\,\% más chico de lo
que la tabla sugiere.

\begin{figure}[!htb]
  \centering
  \includegraphics[width=\textwidth]{fig19_sesgo_del_denominador.pdf}
  \caption{Validación del reparto intraprovincial contra los nacimientos
  departamentales reales del RENAPER. El error no es aleatorio: crece a medida
  que el departamento se achica, y en la dirección que deprime la tasa de los
  lugares chicos.}
  \label{fig:sesgo}
\end{figure}

Conviene además ser explícito sobre lo que \emph{no} se pudo verificar. La serie
del RENAPER parece estar construida por residencia o registro y no por lugar de
ocurrencia: su tasa bruta de natalidad departamental tiene media 15,5 por mil y
el 97\,\% de los casos entre 8 y 30, valores incompatibles con un conteo por
lugar del parto. Validar un reparto por población residente contra una fuente por
residencia es en parte circular. A nivel provincial las dos series coinciden
dentro del 1,3\,\% (CABA: DEIS \num{397191}, RENAPER \num{402492}), de modo que
con los datos disponibles \textbf{no se puede verificar} la afirmación de que
ambas puntas del cociente usan la misma definición.

Esto importa porque el portal del DEIS publica la serie como nacimientos
\emph{ocurridos}, por lugar del parto, que sería la misma definición que usa el
\texttt{P19} de Wikidata: quien nació en una maternidad de la Capital figuraría
en la Capital en las dos puntas del cociente y el cociente sería válido. Pero esa
simetría, aun si se cumpliera, \textbf{solo vale a nivel provincial}, que es donde
el DEIS es dato real. Por debajo de la provincia el denominador ya no es por
ocurrencia: es población residente repartida. Y es justamente ahí —departamento,
ciudad, tamaño de localidad— donde vive H1. La defensa metodológica más fuerte del
estudio no aplica en el nivel donde se mide el efecto principal; se declara como
la limitación central (§\ref{sec:limitaciones}).

Como control se reportan además cinco denominadores alternativos: población total
en los censos 1991, 2001, 2010 y 2022, y personas que en 2022 declararon haber
nacido en cada provincia (variable \texttt{P14}). \textbf{El orden de las
jurisdicciones no cambia con ninguno de ellos} (figura~\ref{fig:sensibilidad}).

\begin{figure}[!htb]
  \centering
  \includegraphics[width=0.92\textwidth]{fig24_sensibilidad_denominador.pdf}
  \caption{Producción observada sobre esperada con seis denominadores distintos.
  Si el resultado dependiera del baseline elegido, las jurisdicciones se
  cruzarían de posición; no lo hacen.}
  \label{fig:sensibilidad}
\end{figure}

\subsection{Ventana de cohortes}

\textbf{1975--2008.} Los dos límites los impuso el dato. Hacia atrás, la serie
del DEIS no publica 1971--1974, único hueco en 110 años: empezar en 1970 dejaba a
esa cohorte con el denominador de un solo año y su tasa salía cinco veces más
alta sin que nada fallara. El pipeline ahora rechaza cualquier ventana con
huecos. Hacia adelante, la carrera: quien nació en 2008 tenía 14 años en 2022 y
no puede haber debutado.

La censura a derecha se cuantifica cohorte por cohorte
(figura~\ref{fig:censura}): la tasa sube de 17 por \num{100000} en 1975--1979 a
un pico de 37 en 1985--1989, y cae a 15 en 2000--2004 y 3 en 2005--2009. La caída
final es censura; la inicial es cobertura, porque Wikidata registra peor a los
jugadores más viejos. El análisis principal usa toda la ventana y se repite
restringido a cohortes $\leq 2002$ como robustez.

\begin{figure}[!htb]
  \centering
  \includegraphics[width=0.88\textwidth]{fig12_censura_por_cohorte.pdf}
  \caption{Diagnóstico de censura y cobertura por quinquenio de nacimiento. Las
  dos caídas tienen causas distintas y solo una es un artefacto de diseño.}
  \label{fig:censura}
\end{figure}

\subsection{Muestra}

Wikidata vía SPARQL, paginado por año de nacimiento, con snapshot fechado.

\begin{center}\small
\begin{tabular}{@{}lr@{}}
\toprule
\textbf{Paso} & \textbf{$n$}\\
\midrule
Futbolistas argentinos en Wikidata & \num{9115}\\
Con precisión de fecha $\geq$ año & \num{8976}\\
Muestra masculina & \num{8649}\\
Con lugar de nacimiento (\texttt{P19}) & \num{8290}\\
Con lugar resuelto dentro de Argentina & \num{7711}\\
\textbf{Cohortes 1975--2008} & \textbf{\num{5511}}\\
\bottomrule
\end{tabular}
\end{center}

\subsection{Resolución geográfica}

Los nombres de lugares argentinos en Wikidata son inconsistentes: hay homónimos
entre provincias —cuatro localidades se llaman Santa Rosa— y entidades que son
provincias o partidos y no ciudades. \textbf{Se resuelve por coordenada}
(\texttt{P625}) contra la API Georef del Estado argentino; el nombre solo se usa
como chequeo cruzado.

La granularidad de cada entidad se clasifica antes de usarla. El caso más
consecuente es la entidad «Argentina», que 255 jugadores tienen como lugar de
nacimiento: su centroide cae en el departamento Presidente Roque Sáenz Peña de
Córdoba y, sin ese filtro, convertía a General Levalle —\num{5674} habitantes— en
la tercera cuna de futbolistas del país.

El mismo problema aparece un nivel más arriba y también se corrige. Cuando el
\texttt{P19} es una provincia entera, su centroide cae en un departamento
cualquiera: el de Buenos Aires en Azul, el de Córdoba en Tercero Arriba, el de
San Juan en Ullum. Los 110 jugadores en esa situación quedaban asignados a esos
departamentos, que así llegaban al top-12 nacional de tasa —Ullum y Tumbaya con
la totalidad de su cuenta fabricada, Azul inflado un 222\,\%—. Ahora conservan la
provincia, que es dato válido, y quedan excluidos del análisis departamental y
del de tamaño de ciudad. \textbf{El análisis provincial usa \num{5511}
jugadores; el departamental, \num{5389}; el de tamaño de ciudad, \num{5248}.}

\textbf{«Tamaño de ciudad» significa aglomerado urbano cuando existe.} Lanús no
es una ciudad de \num{200000} habitantes: es una porción de un conurbano de 16,2
millones. Se usa la definición de aglomerado del propio INDEC y se reporta la
variante con la localidad aislada como control.

\subsection{Estadística}

Chi-cuadrado de bondad de ajuste contra la distribución real de nacimientos
—nunca uniforme—, con Cohen's $w$ y Cramér's $V$. Tasas por \num{100000} con
intervalo exacto de Poisson (Garwood), elegido porque muchos departamentos tienen
cero, uno o dos jugadores. Razones de tasas y \emph{odds ratios} con IC 95\,\%.
Regresión binomial negativa con \emph{offset} logarítmico de los nacimientos.
Proporciones con intervalo binomial exacto (Clopper-Pearson). Corrección de
Benjamini-Hochberg en los cruces exploratorios.

\subsection{La limitación central: nacer no es formarse}

Wikidata da el lugar de nacimiento, no el de desarrollo. Para H3 se construyó un
proxy del club formador: el vínculo jugador-club (\texttt{P54}) con la fecha de
inicio más temprana, excluyendo selecciones. \textbf{Es un proxy.} Wikidata suele
omitir las inferiores, con lo cual el primer club listado es muchas veces el de
debut profesional; y su cobertura es muy desigual por nivel: 99,2\,\% entre los
jugadores de selección mayor contra 12,7\,\% en el resto. H3 se apoya en
\num{1923} jugadores con origen y club formador ubicados en Argentina.

% =============================================================================
\section{Resultados}

\subsection{H1 — El efecto aparece invertido, y es un escalón}

\begin{table}[!htb]\centering\small
\caption{Producción por tramo de tamaño de la ciudad de nacimiento.
$\chi^2(4) = 452{,}2$; $p < 10^{-95}$; $w = 0{,}29$; $n = \num{5248}$.}
\begin{tabular}{@{}lrrrll@{}}
\toprule
\textbf{Tamaño} & \textbf{Futbolistas} & \textbf{Nacidos} & \textbf{Tasa} &
\textbf{IC 95\,\%} & \textbf{RR vs $>$500k}\\
\midrule
$<$10k    & 431 & \num{3339078}  & 12,9 & 11,7--14,2 & 0,42 (0,38--0,47)\\
10--50k   & 568 & \num{3418174}  & 16,6 & 15,3--18,0 & 0,55 (0,50--0,60)\\
50--100k  & 275 & \num{1280832}  & 21,5 & 19,0--24,2 & 0,70 (0,62--0,80)\\
100--500k & 631 & \num{2934614}  & 21,5 & 19,9--23,3 & 0,71 (0,65--0,77)\\
\textbf{$>$500k} & \textbf{3343} & \num{10975566} & \textbf{30,5} & 29,4--31,5 & 1,00\\
\bottomrule
\end{tabular}
\end{table}

El tramo de \num{50000} a \num{100000} habitantes —el óptimo que predice la
literatura— no muestra ningún pico (figura~\ref{fig:tramos}). Con los cortes de
Côté et al.\ el resultado es el mismo: de 8,4 por \num{100000} en localidades de
menos de \num{1000} habitantes a 30,5 en las de más de \num{500000}.

\begin{figure}[!htb]
  \centering
  \includegraphics[width=0.88\textwidth]{fig04_tasa_por_tramo.pdf}
  \caption{Tasa de producción por tramo de tamaño, con intervalos exactos de
  Poisson. El óptimo que predice la literatura internacional no aparece.}
  \label{fig:tramos}
\end{figure}

\textbf{No es un gradiente: es un escalón.} La tabla de cinco tramos sugiere una
progresión ordenada, pero esa forma la produce el ancho de las categorías. Por
decil de tamaño de ciudad las tasas son 22,6 \textbullet{} 4,6 \textbullet{} 14,0
\textbullet{} 4,9 \textbullet{} 8,5 \textbullet{} 8,8 \textbullet{} 13,3
\textbullet{} 11,3 \textbullet{} 16,0 \textbullet{} 25,9: tres de los nueve pasos
bajan, y los nueve deciles por debajo de \num{10000} habitantes —463 futbolistas,
el 9\,\% de la muestra— no muestran ninguna tendencia (figura~\ref{fig:deciles}).
Todo el efecto es el salto del decil superior.

Conviene registrarlo porque \textbf{un escalón único es también la forma que
produciría el registro del parto en la ciudad cabecera}
(§\ref{sec:limitaciones}), mientras que un gradiente continuo no lo sería.

\begin{figure}[!htb]
  \centering
  \includegraphics[width=\textwidth]{fig20_deciles_sin_gradiente.pdf}
  \caption{La misma relación desagregada en diez deciles en lugar de cinco
  tramos. La progresión ordenada de la tabla anterior no sobrevive a la
  desagregación.}
  \label{fig:deciles}
\end{figure}

\textbf{Regresión.} Modelo binomial negativo sobre \num{3459} ciudades: cada
\emph{e-fold} de tamaño multiplica la tasa por 1,181 (\ic{1,120}{1,246};
$p < 0{,}0001$). Con el término cuadrático ni el lineal ni el cuadrático resultan
significativos ($p = 0{,}11$ y $p = 0{,}44$) y el AIC empeora (3654,3 contra
3652,8). No hay curva en U invertida que ajustar.

\textbf{Y el tamaño del efecto es chico.} El pseudo-$R^2$ de McFadden de ese
modelo es \textbf{0,011} y la devianza explicada, 2,6\,\%: el tamaño de la ciudad
da cuenta de alrededor del 1\,\% de la variación entre ciudades
(figura~\ref{fig:scatter}). La pendiente es real y está bien estimada; su
capacidad para predecir cuántos futbolistas produce una ciudad determinada es
casi nula. A \num{100000} habitantes las ciudades reales van de 3 a 100 por cada
\num{100000} nacidos.

\begin{figure}[!htb]
  \centering
  \includegraphics[width=\textwidth]{fig07_scatter_tamano_tasa.pdf}
  \caption{Tamaño de la ciudad contra tasa de producción, ciudad por ciudad. La
  dispersión vertical a cada tamaño es mucho mayor que la pendiente.}
  \label{fig:scatter}
\end{figure}

\textbf{Robustez.} El patrón se sostiene al cambiar la unidad de agregación y la
ventana de cohortes: RR 0,42 con aglomerado urbano, 0,39 con localidad censal
aislada, 0,41 restringido a cohortes $\leq 2002$. El hallazgo no es una decisión
de agregación.

\subsection{H2 — AMBA y corredor pampeano}

\begin{table}[!htb]\centering\small
\caption{Producción por región. $\chi^2(5) = 1050{,}7$; $w = 0{,}44$.}
\begin{tabular}{@{}lrrll@{}}
\toprule
\textbf{Región} & \textbf{Futbolistas} & \textbf{Tasa} & \textbf{IC 95\,\%} &
\textbf{RR vs AMBA}\\
\midrule
AMBA      & 1859 & 35,0 & 33,5--36,7 & 1,00\\
Pampeana  & 2600 & 29,5 & 28,4--30,7 & 0,84\\
Cuyo      &  265 & 15,2 & 13,4--17,2 & 0,43\\
Patagonia &  153 & 12,7 & 10,8--14,9 & 0,36\\
NEA       &  255 &  9,6 &  8,4--10,8 & 0,27\\
NOA       &  269 &  8,2 &  7,2-- 9,2 & 0,23\\
\bottomrule
\end{tabular}
\end{table}

\textbf{H2 no se sostiene.} El interior no produce más que el AMBA: produce
menos, y en el norte produce cuatro veces menos. El mapa departamental
(figura~\ref{fig:mapa}) muestra el corredor con nitidez, y el orden provincial
(figura~\ref{fig:ranking}) es estable con cualquiera de los seis denominadores
probados.

\begin{figure}[!htb]
  \centering
  \includegraphics[height=0.83\textheight]{fig01_mapa_departamentos_tasa.pdf}
  \caption{Producción per cápita por departamento. El patrón dominante es un
  corredor AMBA--Rosario--Santa Fe, no una dispersión federal.}
  \label{fig:mapa}
\end{figure}

\begin{figure}[!htb]
  \centering
  \includegraphics[width=0.84\textwidth]{fig23_ranking_provincias.pdf}
  \caption{Las 24 jurisdicciones ordenadas por tasa, con intervalos exactos de
  Poisson. A nivel provincial el denominador es dato real del DEIS, sin
  estimación: es el nivel más sólido del estudio.}
  \label{fig:ranking}
\end{figure}

\textbf{Una advertencia sobre el ranking departamental.} No se aplica
\emph{shrinkage} ni Bayes empírico, de modo que un departamento con dos jugadores
encabeza cualquier ranking per cápita por puro ruido de Poisson. La
figura~\ref{fig:funnel} separa qué departamentos se apartan de verdad de la media
nacional y cuáles son varianza; el ranking crudo no debe leerse sin ella.

\begin{figure}[!htb]
  \centering
  \includegraphics[width=\textwidth]{fig21_funnel_departamentos.pdf}
  \caption{\emph{Funnel plot} de los departamentos. Dentro del embudo, la tasa es
  compatible con la media nacional por más extrema que parezca en un ranking.}
  \label{fig:funnel}
\end{figure}

\subsection{H3 — La formación está mucho más concentrada que el nacimiento}

Comparado con la población general, el futbolista migra cinco veces más: 47,1\,\%
se forma fuera de su provincia de nacimiento contra 13,8\,\% de la población
general que reside fuera de su provincia de nacimiento (OR 5,58;
\ic{5,10}{6,10}). El corte por origen es un escalón, no un gradiente: quien nace
en un aglomerado de más de \num{500000} habitantes se forma a 7\,km de mediana;
quien nace en cualquier otro lado cambia de departamento en el 86--98\,\% de los
casos.

\begin{figure}[!htb]
  \centering
  \includegraphics[height=0.80\textheight]{fig08_flujos_nacimiento_club.pdf}
  \caption{Flujos entre la provincia de nacimiento y la del club formador. El
  tamaño del nodo es el saldo neto entre los que recibe y los que pierde.}
  \label{fig:flujos}
\end{figure}

\begin{figure}[!htb]
  \centering
  \includegraphics[width=0.86\textwidth]{fig25_matriz_flujo_regiones.pdf}
  \caption{Matriz de flujo entre regiones. La diagonal es la retención: el AMBA
  se queda con el 92\,\% de los suyos y el NEA con el 9\,\%.}
  \label{fig:matriz}
\end{figure}

\textbf{Esta subsección tiene un problema de muestra que hay que declarar antes
de leer los números.} La cobertura del club formador va del 99,2\,\% en jugadores
de selección al 12,7\,\% en el resto, de modo que la submuestra de H3 está
enriquecida al doble en jugadores de elite y vaciada cinco veces del resto —
precisamente los que se fueron a los clubes grandes. El 47,1\,\% de migración y
la concentración en unos pocos clubes describen a los que llegaron lejos, no a la
población de futbolistas. Deben leerse como orden de magnitud.

\subsection{H4 — No lo fabrica la cobertura de Wikidata}

La objeción más seria contra todo lo anterior es el sesgo de cobertura. La prueba
está en los jugadores de selección mayor, donde el registro de Wikidata es
prácticamente censal.

El patrón por tamaño se sostiene en los cuatro niveles competitivos, incluida la
selección ($\chi^2(4) = 26{,}5$; $p < 0{,}0001$; $w = 0{,}32$; $n = 255$): 15,3
seleccionados por millón de nacidos en aglomerados de más de \num{500000} contra
7,5 en localidades de menos de \num{10000} (figura~\ref{fig:seleccion}). Por
región, 17 por millón en el AMBA contra 2 en el NOA. En cambio \textbf{H4 tal
como estaba formulada no se sostiene}: no hay evidencia de que la elite provenga
más de ciudades chicas (OR 0,89; \ic{0,71}{1,13}).

\begin{figure}[!htb]
  \centering
  \includegraphics[width=\textwidth]{fig13_seleccion.pdf}
  \caption{Producción de jugadores de selección mayor por tamaño de ciudad y por
  región. Es el control del sesgo de cobertura del corpus.}
  \label{fig:seleccion}
\end{figure}

Un matiz que la tabla de cinco tramos esconde: dentro de la selección, los cuatro
tramos no metropolitanos son indistinguibles entre sí (7,5 \textbullet{} 7,0
\textbullet{} 8,6 \textbullet{} 9,2 por millón, con intervalos que se solapan por
completo). Como en H1, lo que hay es un contraste binario contra los grandes
aglomerados, no una escalera.

Vale una aclaración sobre el alcance de esta prueba: acota \textbf{la amenaza de
cobertura} y no dice nada sobre la otra —el registro del parto en la ciudad
cabecera—, que afecta a los seleccionados exactamente igual que al resto.

\subsection{De los juveniles a la Mayor: el resultado que no depende del denominador}
\label{sec:conversion}

Todo lo anterior son tasas por nacido y por lo tanto heredan los dos problemas
del denominador. Este análisis no.

La pregunta es otra: entre los futbolistas que ya llegaron a un juvenil de la
selección (sub-17 o sub-20), ¿qué proporción llega después a la Mayor, según
dónde nacieron? El denominador acá no es una estimación de nacimientos: son los
347 juveniles observados.

\begin{table}[!htb]\centering\small
\caption{Conversión de juvenil a Selección Mayor según el tamaño de la ciudad de
nacimiento. Intervalos binomiales exactos.}
\begin{tabular}{@{}lrrrl@{}}
\toprule
\textbf{Ciudad de nacimiento} & \textbf{Juveniles} & \textbf{Llegan a la Mayor} &
\textbf{\%} & \textbf{IC 95\,\%}\\
\midrule
$<$10k    &  19 &  7 & 36,8 & 16,3--61,6\\
10--50k   &  25 & 14 & 56,0 & 34,9--75,6\\
50--100k  &  20 &  6 & 30,0 & 11,9--54,3\\
100--500k &  31 & 12 & 38,7 & 21,8--57,8\\
\textbf{$>$500k} & \textbf{242} & \textbf{68} & \textbf{28,1} & 22,5--34,2\\
\bottomrule
\end{tabular}
\end{table}

Agrupando en el contraste que tiene potencia —fuera de un gran aglomerado contra
adentro—: \textbf{41,9\,\% (44 de 105) contra 28,1\,\% (68 de 242)}. OR 1,85
(\ic{1,14}{2,98}); $\chi^2(1) = 5{,}77$; $p = 0{,}016$; Fisher exacto
$p = 0{,}013$ (figura~\ref{fig:conversion}).

\textbf{No es un efecto generacional.} Una regresión logística que agrega el año
de nacimiento como control deja el OR en 1,85 (\ic{1,13}{3,02}; $p = 0{,}014$).

\begin{figure}[!htb]
  \centering
  \includegraphics[width=\textwidth]{fig14_conversion_juvenil_mayor.pdf}
  \caption{Conversión de juvenil a Selección Mayor por tamaño de la ciudad de
  nacimiento. El denominador son jugadores observados, no nacimientos estimados.}
  \label{fig:conversion}
\end{figure}

Por qué este resultado resiste lo que los otros no:

\begin{itemize}[leftmargin=1.4em,itemsep=0.2em]
  \item \textbf{No usa denominador poblacional.} El sesgo de imputación del
    17\,\% no lo toca.
  \item \textbf{No depende de dónde ocurrió el parto} para construir una tasa: el
    registro del nacimiento en la cabecera clasificaría a un jugador de pueblo
    como $>$500k, y eso \emph{atenúa} el contraste en vez de crearlo. El efecto
    medido es un piso.
  \item \textbf{La cobertura de Wikidata es del 97\,\%} entre jugadores de
    selección, y es la misma para los dos grupos que se comparan.
  \item \textbf{Compara dentro de un grupo ya filtrado por talento reconocido},
    de modo que no hay que suponer que el talento latente se distribuye parejo
    entre regiones.
\end{itemize}

La figura~\ref{fig:embudo} pone los tres escalones juntos y muestra dónde se
invierte el signo: en los dos primeros —que dependen del denominador— el AMBA
gana; en el tercero —que no— pierde.

\begin{figure}[!htb]
  \centering
  \includegraphics[width=\textwidth]{fig16_embudo_por_origen.pdf}
  \caption{El embudo completo, del nacimiento a la Selección Mayor. Los dos
  primeros paneles son tasas por nacido; el tercero es una proporción dentro de
  un grupo observado.}
  \label{fig:embudo}
\end{figure}

\subsection{Exploratorio: posición y región}

\textbf{Estrictamente exploratorio.} De 24 contrastes, seis sobreviven a la
corrección de Benjamini-Hochberg, y todos involucran a las dos regiones con más
casos: el AMBA produce más defensores y menos mediocampistas de lo esperado, y la
región pampeana lo inverso. \textbf{El mito de «las delanteras del norte» no
aparece}: ningún contraste que involucre al NOA o al NEA sobrevive a la
corrección.

% =============================================================================
\section{Discusión}

\subsection{Por qué el efecto está invertido}

El \emph{birthplace effect} clásico se apoya en un supuesto implícito: que la
infraestructura de desarrollo deportivo está razonablemente distribuida, de modo
que lo que diferencia a los lugares es la calidad del entorno de juego informal.
Bajo ese supuesto la ciudad mediana gana: tiene espacio y tiene liga.

En Argentina ese supuesto no se cumple. La formación está concentrada: diez
clubes concentran la mitad de la formación registrada del país, y están todos en
el AMBA, Rosario o La Plata. La interpretación natural es que el lugar de
nacimiento no mide acá la calidad del entorno formativo sino \textbf{la distancia
a la infraestructura formativa}, y que esa distancia opera como filtro de acceso.

Es una interpretación, no un resultado, y conviene ser explícito. El estudio
\textbf{no mide distancia a un club con inferiores}: ninguna regresión incluye esa
variable, ni nivel socioeconómico, ni densidad de ligas locales, ni existencia de
pensión. El análisis es descriptivo y las asociaciones no están controladas.

El propio mecanismo permite dos pruebas, y dan resultados distintos:

\begin{itemize}[leftmargin=1.4em,itemsep=0.25em]
  \item \emph{Si la causa fuera la centralización creciente de las inferiores, la
    brecha debería ensancharse por cohorte.} No lo hace
    (figura~\ref{fig:cohorte}): la razón de tasas entre pueblo y gran aglomerado
    es 0,34 \textbullet{} 0,43 \textbullet{} 0,41 \textbullet{} 0,55 en las
    cohortes de los 70, 80, 90 y 2000, plana y con intervalos solapados. El
    mecanismo no queda refutado —la centralización pudo ser estable en el
    período— pero tampoco encuentra apoyo donde debería.
  \item \emph{Si el acceso filtrara por algo distinto del talento, quien
    atraviesa el filtro más exigente debería rendir mejor.} Eso sí aparece, y es
    el resultado más robusto del trabajo (§\ref{sec:conversion}).
\end{itemize}

\begin{figure}[!htb]
  \centering
  \includegraphics[width=\textwidth]{fig22_efecto_por_cohorte.pdf}
  \caption{Evolución del efecto por década de nacimiento. La centralización
  creciente de las inferiores predice una brecha que se ensancha; la serie está
  plana.}
  \label{fig:cohorte}
\end{figure}

Esa asimetría —menos acceso, mejor rendimiento condicional— es la evidencia más
directa de que el patrón geográfico principal refleja acceso y no distribución de
talento. Y, a diferencia de todo lo demás, no depende del denominador.

Frente a la literatura, el aporte hay que dimensionarlo con cuidado. La revisión
sistemática de Hernández-Simal y colegas \cite{revision2024} ya documenta
hallazgos que favorecen a las áreas densas, ya reporta que en fútbol no hay
ventaja consistente de las ciudades chicas, y ya identifica la proximidad a
centros de rendimiento como uno de sus tres ejes explicativos. Este trabajo
\textbf{no descubre un mecanismo nuevo}: aporta la primera medición argentina
—Argentina no aparece en esa revisión—, con un denominador de nacidos vivos por
cohorte que es mejor que el de buena parte de la literatura, y un diseño
condicional que no necesita denominador.

\subsection{Qué cambió al corregir el denominador}

Vale registrarlo porque es el tipo de cosa que decide un resultado en silencio.
Con la población censada en 2022 como denominador, el RR de las localidades de
menos de \num{10000} habitantes contra los grandes aglomerados daba 0,57; con los
nacimientos reales da 0,42. La diferencia es la migración: los pueblos son
exportadores netos de población, así que contar a sus residentes de 2022
subestimaba cuánta gente había nacido ahí e inflaba su tasa. En la misma
corrección, el AMBA pasó de estar por debajo de la región pampeana a estar por
encima: las áreas metropolitanas tienen menos hijos por habitante.

\subsection{Limitaciones}
\label{sec:limitaciones}

\begin{enumerate}[leftmargin=1.6em,itemsep=0.3em]
  \item \textbf{Nacer no es criarse, y en Argentina el parto ocurre donde hay
    maternidad.} Es la limitación más seria y no está acotada. Un chico de un
    pueblo de \num{3000} habitantes nace en la cabecera departamental y queda
    registrado ahí, en el DEIS y en Wikidata. Eso vacía sistemáticamente a las
    localidades chicas y llena a las cabeceras. Dos señales de que el problema es
    real: el tramo $<$10k se construye sobre un padrón que incluye 264
    «localidades» de menos de 100 habitantes —entre ellas parajes de seis
    personas, donde materialmente no nace nadie— y la forma del efecto es un
    escalón único, que es exactamente lo que este artefacto produciría.
    \textbf{El análisis de §\ref{sec:conversion} se diseñó para esquivar el
    problema; los demás lo padecen.}
  \item \textbf{El reparto intraprovincial es un supuesto con sesgo
    % El «%» de babel-spanish mide en pt y `\,` en modo matemático usa mu:
    % mezclarlos da «Incompatible glue units». El signo va en matemático y el
    % símbolo afuera.
    direccional}: $+17$\,\% en el decil de departamentos más chicos y nulo en el
    más grande (§\ref{sec:denominador}), empujando en la misma dirección que el
    hallazgo.
  \item \textbf{Cobertura de Wikidata.} Es un corpus de notabilidad, no un
    registro, y nunca se contrastó contra un padrón independiente de futbolistas
    profesionales: \textbf{la tasa de error del \texttt{P19} no está medida}. La
    cobertura además varía al doble entre cohortes.
  \item \textbf{H3 se mide sobre una muestra seleccionada por el desenlace} (ver
    §3.3). Sus números son orden de magnitud, no estimaciones.
  \item \textbf{La comparación de H3 con la población general no es estrictamente
    comparable}: 47,1\,\% (futbolistas, nacimiento a primer club, alrededor de
    los 18 años) contra 13,8\,\% (toda la población, todas las edades, nacimiento
    a residencia 2022). La variable \texttt{P14} del censo no está cruzada con
    edad.
  \item \textbf{Sin controles.} No hay ninguna covariable más allá del tamaño de
    la ciudad. Las lecturas causales de esta sección son interpretaciones, no
    estimaciones.
  \item \textbf{Tamaño de efecto chico} donde se lo mide de forma continua:
    pseudo-$R^2 = 0{,}011$.
  \item \textbf{Números chicos en el mapa departamental}: sin \emph{shrinkage}
    (ver figura~\ref{fig:funnel}).
  \item \textbf{Censura a derecha} en las cohortes 2003--2008, cuantificada y
    repetida sin ellas.
  \item \textbf{Unidad geográfica en metrópolis fragmentadas}: las tasas por
    departamento se inflan en el núcleo de los aglomerados que cruzan límites
    administrativos.
  \item \textbf{Solo fútbol masculino.}
\end{enumerate}

\subsection{Implicancias}

Conviene separar lo que el dato sostiene de lo que sugiere.

\textbf{Lo que el dato sostiene.} Entre los futbolistas que ya llegaron a un
juvenil de la selección, los nacidos fuera de un gran aglomerado llegan a la
Mayor con más frecuencia (OR 1,85; $p = 0{,}013$; sin cambios al controlar por
cohorte). Es un contraste dentro de un grupo observado, sin denominador estimado.
Si el acceso midiera talento, la tasa de conversión debería ser igual en los dos
grupos; no lo es. \textbf{Eso es evidencia de que el filtro de acceso está
dejando afuera jugadores que habrían rendido.}

Para quien tiene que decidir dónde poner un centro de detección, la implicancia
es concreta y no depende de ninguna de las limitaciones anteriores: \textbf{un
juvenil del interior es, en promedio, mejor apuesta que uno del AMBA con el mismo
nivel alcanzado.} Con 105 casos fuera del AMBA el intervalo es ancho
(1,14--2,98) y merece confirmarse con el padrón real de convocatorias de AFA, que
existe y no es público.

\textbf{Lo que el dato solo sugiere.} Que las regiones que producen un cuarto de
lo esperado sean «reservas desaprovechadas» requiere suponer que el talento
latente se distribuye parejo entre regiones —un supuesto razonable pero no
testeado acá—. La dirección del argumento es plausible; su magnitud, no está
establecida.

Y una advertencia metodológica que sí se sostiene sola: \textbf{cualquier métrica
de producción basada en el lugar de nacimiento le atribuye al AMBA jugadores que
el AMBA absorbió, no formó} y, en la medida en que el parto ocurre en la
cabecera, también le atribuye a las ciudades chicos que nacieron ahí de
casualidad.

\subsection{Qué haría falta para cerrar el argumento}

Dos cosas, en este orden:

\begin{enumerate}[leftmargin=1.6em,itemsep=0.2em]
  \item \textbf{Medir la tasa de error del \texttt{P19}} contra una fuente
    independiente, sobre una muestra estratificada por tamaño de ciudad. Sin eso,
    todo lo que depende del lugar de nacimiento descansa sobre un dato sin
    validar.
  \item \textbf{Acotar el artefacto de maternidad}, contrastando nacimientos por
    ocurrencia contra nacimientos por residencia de la madre en el nivel
    geográfico más fino que publique la fuente.
\end{enumerate}

El resultado de §\ref{sec:conversion} no depende de ninguna de las dos.

% =============================================================================
\section{Reproducibilidad}

Todo el pipeline es código; \texttt{data/raw/} no se edita y cada descarga deja un
manifiesto fechado con URL y SHA-256. Las decisiones de diseño están en
\texttt{config.yaml} con su justificación, y las funciones que comparten
numerador y denominador tienen tests. El orden de ejecución está en el README.
Las 26 figuras se generan con dos módulos y se exportan en PDF, SVG y PNG.

\paragraph{Fuentes.}
Wikidata vía SPARQL (CC0, snapshot 2026-07-30);
DEIS, serie histórica de nacidos vivos por jurisdicción 1914--2024;
RENAPER, nacimientos por departamento 2012--2022 (validación);
INDEC, Censo 2022 y radios censales 1991--2022;
API Georef e IGN.
Transfermarkt no se utilizó: sus términos prohíben el scraping automatizado.

\begin{thebibliography}{9}
\small
\bibitem{cote2006}
Côté, J., MacDonald, D. J., Baker, J. \& Abernethy, B. (2006).
When «where» is more important than «when»: birthplace and birthdate effects on
the achievement of sporting expertise.
\emph{Journal of Sports Sciences}, 24(10), 1065--1073.

\bibitem{revision2024}
Hernández-Simal, L., Calleja-González, J., Lorenzo Calvo, A. \&
Aurrekoetxea-Casaus, M. (2024).
Birthplace Effect in Soccer: A Systematic Review.
\emph{Journal of Human Kinetics}.

\bibitem{geografia2022}
\emph{The geography of talent development} (2022).
\emph{Frontiers in Sports and Active Living}. PMC9582327.
\textcolor{gray}{[autoría y año verificados solo parcialmente]}
\end{thebibliography}

\end{document}
